\documentclass[11pt,a4paper]{article}

% ──────────────────────────────────────────────────────────────────────────
% AXON-RFC-006 Appendix A — easynet.pages
% Status: Draft v0.1 (technical design + implementation plan)
% Build:  xelatex AXON-RFC-006-A-easynet-pages.tex
% Companion: AXON-RFC-006-stateful-easynet.tex (minimal normative core)
% ──────────────────────────────────────────────────────────────────────────

\usepackage[margin=1in]{geometry}
\usepackage{amsmath,amssymb,amsthm}
\usepackage{xcolor}
\usepackage[colorlinks=true,linkcolor=blue!60!black,citecolor=blue!60!black,urlcolor=blue!60!black]{hyperref}
\usepackage{booktabs}
\usepackage{longtable}
\usepackage{array}
\usepackage{tabularx}
\usepackage{xltabular}
\newcolumntype{L}[1]{>{\raggedright\arraybackslash}p{#1}}
\usepackage{listings}
\usepackage{enumitem}
\usepackage{tcolorbox}
\usepackage{titlesec}
\usepackage{fancyhdr}
\usepackage{parskip}
\usepackage{xspace}

\usepackage{fontspec}
\usepackage{xeCJK}
\setCJKmainfont{PingFang SC}[
    BoldFont={PingFang SC Semibold},
    ItalicFont={PingFang SC Light}
]

\pagestyle{fancy}
\fancyhf{}
\fancyhead[L]{\small AXON-RFC-006 Appendix A}
\fancyhead[R]{\small easynet.pages --- Technical Design}
\fancyfoot[C]{\thepage}
\renewcommand{\headrulewidth}{0.4pt}

\titleformat{\section}{\Large\bfseries\color{black!85}}{\thesection}{0.7em}{}
\titleformat{\subsection}{\large\bfseries\color{black!75}}{\thesubsection}{0.6em}{}
\titleformat{\subsubsection}{\normalsize\bfseries\color{black!70}}{\thesubsubsection}{0.5em}{}

\definecolor{normativecolor}{RGB}{180,30,30}
\definecolor{derivedcolor}{RGB}{30,90,180}
\definecolor{deferredcolor}{RGB}{120,90,30}
\definecolor{codebg}{RGB}{245,245,248}
\definecolor{codeborder}{RGB}{210,210,220}

\newcommand{\normsv}{\textcolor{normativecolor}{\textbf{[normative]}}\xspace}
\newcommand{\derived}{\textcolor{derivedcolor}{\textbf{[derived]}}\xspace}
\newcommand{\deferred}{\textcolor{deferredcolor}{\textbf{[deferred]}}\xspace}

\lstdefinestyle{ealstyle}{
    backgroundcolor=\color{codebg},
    basicstyle=\ttfamily\small,
    frame=single,
    rulecolor=\color{codeborder},
    breaklines=true,
    showstringspaces=false,
    keywordstyle=\color{blue!70!black}\bfseries,
    commentstyle=\color{green!40!black}\itshape,
    stringstyle=\color{red!60!black},
    morekeywords={state_type,state_key,canonical,operational,transition,receipt,view,principal,agent,owner_signature,attempt_id,transition_id,fn,let,struct,impl,pub,async,await},
    columns=fullflexible,
    keepspaces=true,
    aboveskip=0.6em,
    belowskip=0.6em,
}
\lstset{style=ealstyle}

\newtheorem{theorem}{Theorem}
\newtheorem{lemma}[theorem]{Lemma}
\newtheorem{corollary}[theorem]{Corollary}
\theoremstyle{definition}
\newtheorem{definition}{Definition}

\newtcolorbox{invariant}[1][Invariant]{
    colback=red!4!white,
    colframe=red!50!black,
    fonttitle=\bfseries,
    title={Invariant --- #1},
    sharp corners=south,
    boxrule=0.6pt,
}
\newtcolorbox{notebox}[1][]{
    colback=blue!4!white,
    colframe=blue!50!black,
    fonttitle=\bfseries,
    title={Note},
    sharp corners=south,
    boxrule=0.6pt,
    #1
}
\newtcolorbox{warningbox}[1][]{
    colback=orange!6!white,
    colframe=orange!70!black,
    fonttitle=\bfseries,
    title={Transitional misalignment},
    sharp corners=south,
    boxrule=0.6pt,
    #1
}
\newtcolorbox{decisionbox}[1][]{
    colback=purple!4!white,
    colframe=purple!60!black,
    fonttitle=\bfseries,
    title={Open decision},
    sharp corners=south,
    boxrule=0.6pt,
    #1
}

\title{
    \vspace{-2em}
    \textbf{\Large AXON-RFC-006 Appendix A: easynet.pages}\\[0.3em]
    \large Technical design, algorithms, and implementation plan
    for the first reference state machine of stateful EasyNet
}
\author{
    Silan Hu \\ \small National University of Singapore
}
\date{Draft v0.1, \today}

\begin{document}
\maketitle

\begin{abstract}
\noindent
This appendix specifies \texttt{easynet.pages}, the simplest
\texttt{state\_type} that exercises the full set of normative
primitives defined in AXON-RFC-006 main. A page is a
\emph{dual-view state object}: it carries identity, canonical
content, and version under the rules of \S\S1--7 of the main
body, and exposes that state simultaneously through
\texttt{human\_view} (an HTML rendering for browsers) and
\texttt{agent\_view} (a structured, machine-actionable rendering
that includes ability surface, EAL hint, and receipt-history
pointer). The appendix covers four layers of detail. First, the
state machine: the field set, the four-state space, and the seven
transitions. Second, the wire schemas of every transition's
arguments and receipt body. Third, the algorithms: canonical
projection, JCS canonicalisation, hash, ETag, optimistic
concurrency, and view projection, each with a worked byte-level
example. Fourth, the implementation plan: a module map for
\texttt{EasyNet-Cli}, a five-PR slice plan (P-A1 through P-A5),
and a complete end-to-end worked example following one page from
its first \texttt{pages.create} through a Hub-rendered HTML
response and an agent-readable URA fetch. Appendix B (the
\texttt{easynet.web\_app} pressure test) and the main body's
thirteen \texttt{TR-INV-*} invariants are referenced throughout
but not redefined; this document is a \emph{reference
implementation specification}, not a re-derivation of the
protocol.
\end{abstract}

\tableofcontents
\vspace{1em}
\hrule
\vspace{1em}

% ──────────────────────────────────────────────────────────────────────────
\section{Goals and non-goals}
\label{sec:goals}

\subsection{Goals}

\normsv{} This appendix specifies an implementation of
\texttt{easynet.pages} that:

\begin{enumerate}[leftmargin=2.2em,itemsep=0.25em]
    \item Conforms to every \normsv{} clause of AXON-RFC-006
    main (\S\S1--7) and discharges every applicable
    \texttt{TR-INV-*} from the companion Appendix~B.
    \item Provides four canonical transitions (\texttt{create},
    \texttt{publish}, \texttt{update}, \texttt{archive},
    \texttt{delete}) and two query abilities (\texttt{get},
    \texttt{list}).
    \item Stores page content as a content-addressed blob,
    keyed by SHA-256, never embedded in canonical state
    (CSH-INV-2).
    \item Exposes both \texttt{human\_view} (HTML rendered from
    the content blob) and \texttt{agent\_view} (a typed JSON
    document containing canonical fields, ability surface, EAL
    hint, and receipt history pointer) for every published
    page.
    \item Is implementable end-to-end in five PR-sized slices
    (\S\ref{sec:impl-plan}).
\end{enumerate}

\subsection{Non-goals}

\derived{} The following are deliberately out of scope:

\begin{itemize}[leftmargin=2.2em,itemsep=0.2em]
    \item \textbf{webapp} (the long-running, runtime-overlay
    pressure test). Covered by Appendix~B.
    \item \textbf{cross-realm or cross-node hubs}. Phase~1 ships
    a single-node Hub that reverse-proxies pages owned by the
    same daemon. Multi-realm hubs are deferred to RFC-006-A v2.
    \item \textbf{authenticated human callers}
    (\texttt{principal/human-auth}). Phase~1 supports
    \texttt{principal/human-anon} only.
    \item \textbf{HTML sanitisation pipelines and CSP policies
    beyond a fixed default}. Owners author HTML; the Hub serves
    it under a single hard-coded CSP. A sanitisation/template
    layer is RFC-006-A v2.
    \item \textbf{TLS termination, rate limiting, multi-instance
    Hub fleets}. The Hub is an axum process binding a single
    port; deployment-grade hardening is operations work, not
    protocol work.
\end{itemize}

% ──────────────────────────────────────────────────────────────────────────
\section{Relation to the normative core}
\label{sec:relation}

This appendix is a \emph{reference implementation} of the
normative core in AXON-RFC-006 main. It does not redefine any
primitive. The mapping is:

\begin{center}
\renewcommand{\arraystretch}{1.2}
\begin{tabularx}{\linewidth}{l X}
\toprule
\textbf{Main clause} & \textbf{This appendix's instantiation} \\
\midrule
\S1 State Object               & \S\ref{sec:state-machine}: pages state\_type, key, lifecycle. \\
\S1.5 SO-INV-1 immutability    & state\_key (realm, owner, slug) is immutable post-create. \\
\S2 Canonical State Hash       & \S\ref{sec:algorithms}: explicit projection $\pi$, JCS, SHA-256 example. \\
\S2 CSH-INV-2 content-as-blob  & \S\ref{sec:storage}: blob layout; canonical state holds only \texttt{content\_hash}. \\
\S3 Receipt Classes            & \S\ref{sec:wire-schemas}: every transition's receipt body shown. \\
\S3 RC-INV-2 optimistic concurrency & \S\ref{sec:algo-concurrency}: \texttt{ETag} = first 16 bytes of canonical hash. \\
\S4 Transition Ability         & \S\ref{sec:wire-schemas}: \texttt{transition\_id}s for every page transition. \\
\S5 Authority + delegation     & Phase 1: owner\_agent signs locally; no delegated execution for pages. \\
\S5 principal/ caller          & \S\ref{sec:hub}: Hub mints \texttt{human-anon} URIs per request. \\
\S6 View Projection            & \S\ref{sec:algo-views}: human\_view and agent\_view algorithms. \\
\S7 Receipt Indexing           & \S\ref{sec:storage}: state\_key index in addition to receipt\_id. \\
\bottomrule
\end{tabularx}
\end{center}

\textbf{Reading discipline.} A statement labelled \normsv{} in this
appendix is binding on the \texttt{pages} implementation only; it
does not propose a new protocol-level rule. Where this document
appears to extend the protocol, it does so by instantiating an
already-\normsv{} clause of the main body.

% ──────────────────────────────────────────────────────────────────────────
\section{State machine}
\label{sec:state-machine}

\subsection{state\_type and state\_key}

\normsv{}

\begin{lstlisting}
state_type:  easynet.page

state_key:
    (realm: string, owner_agent_uri: string, slug: string)

  realm           Logical namespace. Phase 1: a single realm per
                  daemon, named in config (default "default").
  owner_agent_uri Canonical authority. The agent that signs every
                  CanonicalReceipt for this page. Embedded in the
                  key (not a field) per main §1.3.
  slug            Owner-namespaced kebab-case identifier
                  (^[a-z0-9][a-z0-9-]{0,62}[a-z0-9]$).

URA form:
  easynet:///r/<realm>/agent/<owner>/page/<slug>

ABNF (slug):
  slug      = head [ *body tail ]
  head      = ALPHA / DIGIT
  body      = ALPHA / DIGIT / "-"
  tail      = ALPHA / DIGIT
  ALPHA     = %x61-7A         ; lowercase ASCII letters only
\end{lstlisting}

The slug grammar deliberately rejects uppercase, dots, slashes, and
internationalised characters. Two reasons: the slug appears in URLs
with no escape pass; and case-sensitive collisions
(\texttt{Hello} vs \texttt{hello}) are a rich source of admin
confusion. RFC-006-A v2 may relax this.

\subsection{Canonical fields}

\normsv{} Per CSH-INV-2 (main \S2.5), large content is held as a
hash. The canonical state of a page is therefore a small map:

\begin{lstlisting}
canonical_fields(page) = {
    manifest_hash:    bytes(32),    // SHA-256 of the manifest
                                    // sub-document (title,
                                    // description, content_type,
                                    // visibility); see §3.3
    content_hash:     bytes(32),    // SHA-256 of the body bytes,
                                    // null/zero in Draft and after
                                    // archive/delete
    state_value:      string,       // "Draft" | "Published"
                                    //   | "Archived" | "Deleted"
    visibility:       string,       // "private" | "realm" | "public"
    version:          uint64        // monotonic; +1 per
                                    // CanonicalReceipt
}
\end{lstlisting}

The schema annotates every field with \texttt{canonical:true}. The
following \emph{operational} fields exist in Phase 1 but are
\emph{never} part of canonical projection $\pi$
(\S\ref{sec:algorithms}):

\begin{lstlisting}
operational_fields(page) = {
    last_view_at:        timestamp,    // observation only
    view_count_estimate: uint64,       // best-effort
    last_caller_class:   string|null   // "browser" | "bot" | ...
}
\end{lstlisting}

These are computed from \texttt{OperationalReceipt}s emitted by
the Hub during view projection (\S\ref{sec:algo-views}); they are
useful for the agent\_view's audit summary but never enter
\texttt{canonical\_state\_hash}.

\subsection{State space}
\label{sec:state-space}

\begin{center}
\renewcommand{\arraystretch}{1.25}
\begin{tabularx}{\linewidth}{l X}
\toprule
\textbf{state\_value} & \textbf{Meaning} \\
\midrule
\texttt{Draft}       & Created but not yet published. content\_hash may be null or set; not served by Hub. \\
\texttt{Published}   & Live. content\_hash non-null. Served by Hub if visibility allows. \\
\texttt{Archived}    & Previously Published; not served by Hub. content\_hash retained for history. \\
\texttt{Deleted}     & Terminal. content\_hash cleared. No further transitions accepted. \\
\bottomrule
\end{tabularx}
\end{center}

\noindent Allowed canonical transitions (single canonical state
space, no runtime overlay --- pages have no long-running serve
process, unlike \texttt{web\_app}):

\begin{center}
\renewcommand{\arraystretch}{1.2}
\begin{tabularx}{\linewidth}{l l X}
\toprule
\textbf{Transition} & \textbf{Pre $\to$ Post} & \textbf{Notes} \\
\midrule
\texttt{pages.create@v1}   & $\emptyset \to$ Draft     & Mints state key. \\
\texttt{pages.publish@v1}  & Draft $\to$ Published     & Requires non-null content\_hash. \\
\texttt{pages.update@v1}   & Published $\to$ Published & version+1; content\_hash and/or manifest\_hash change. \\
\texttt{pages.archive@v1}  & Published $\to$ Archived  & Reversible by re-publish. \\
\texttt{pages.unpublish@v1}& Archived $\to$ Draft      & Allows resume of editing. \\
\texttt{pages.delete@v1}   & * $\to$ Deleted (terminal)& Clears content\_hash. \\
\bottomrule
\end{tabularx}
\end{center}

\noindent Query abilities (no state mutation, no
\texttt{CanonicalReceipt}, only \texttt{OperationalReceipt} for
audit):

\begin{center}
\renewcommand{\arraystretch}{1.2}
\begin{tabularx}{\linewidth}{l l}
\toprule
\textbf{Ability} & \textbf{Returns} \\
\midrule
\texttt{pages.get@v1}  & full canonical projection + content blob handle \\
\texttt{pages.list@v1} & list of (slug, state\_value, version) for one owner \\
\bottomrule
\end{tabularx}
\end{center}

\begin{invariant}[PA-INV-1]
A page \texttt{Draft} or \texttt{Archived} \textbf{MUST NOT} be
served by any Hub regardless of visibility. The Hub
discriminates on \texttt{state\_value} \emph{before} consulting
visibility; visibility gates only Published pages.
\end{invariant}

\textbf{Why PA-INV-1.} A draft is by definition unfinished; serving
it because visibility happened to be \texttt{public} would surface
work-in-progress as published content. Archived pages are
deliberately un-served --- they exist for history, not access.

\subsection{Lifecycle diagram}

\begin{lstlisting}
                       pages.create
                            |
                            v
            ┌─────────────[Draft]─────────────┐
            |                                 |
   pages.publish                       pages.delete
            |                                 |
            v                                 v
        [Published] ──── pages.update ──> [Published]
            |                                 |
   pages.archive                       pages.delete
            |                                 |
            v                                 v
        [Archived] ── pages.unpublish → [Draft]
            |
   pages.delete
            |
            v
         [Deleted]   (terminal)
\end{lstlisting}

% ──────────────────────────────────────────────────────────────────────────
\section{Wire schemas}
\label{sec:wire-schemas}

Every transition's \texttt{ability} input schema and the
\texttt{receipt\_schema} for its \texttt{CanonicalReceipt} are
listed below. JSON Schema is abbreviated (types, required keys
only); the canonical schemas live in
\texttt{abilities/pages.<verb>.ability.toml}.

\subsection{pages.create@v1}

\begin{lstlisting}
args:
{
  slug:         string (matches §3.1 slug ABNF, required)
  manifest:     {
    title:        string,         required
    description:  string|null,
    content_type: "text/html" | "text/markdown",
                                  required
    visibility:   "private" | "realm" | "public",
                                  required
  }
  initial_body: bytes|null        // utf8 or base64 per content_type;
                                  // null leaves the page empty draft
}

CanonicalReceipt body extends the main-§3.1 base with:
{
  page_state_key:    (realm, owner, slug),
  pre_state_hash:    zero(32),    // creation: pre is empty
  post_state_hash:   sha256(jcs(canonical_fields)),
  state_version:     1,
  manifest_hash:     sha256(jcs(manifest)),
  content_hash:      sha256(initial_body) if initial_body else zero(32),
  blob_handles:      [content_hash] if content_hash != 0 else []
}
\end{lstlisting}

\subsection{pages.publish@v1}

\begin{lstlisting}
args:
{
  slug:    string,
  // optional: replace body before publishing
  body:    bytes|null,
  reason:  string|null   // free-form, recorded in receipt
}

precondition: state_value == Draft AND content_hash != 0
              (or `body` is provided and non-empty)

CanonicalReceipt body:
{
  page_state_key:    ...,
  pre_state_hash:    canonical_state_hash(Draft snapshot),
  post_state_hash:   canonical_state_hash(Published snapshot),
  state_version:     prev+1,
  content_hash:      sha256(new body),
  blob_handles:      [content_hash],
  reason:            args.reason
}
\end{lstlisting}

\subsection{pages.update@v1}

\begin{lstlisting}
args:
{
  slug:        string,
  manifest:    manifest|null,    // partial update permitted;
                                 //   null = unchanged
  body:        bytes|null,       // null = unchanged
  reason:      string|null
}
precondition: state_value == Published AND
              (manifest != null OR body != null)

CanonicalReceipt body:
{
  page_state_key:    ...,
  pre_state_hash:    h_prev,
  post_state_hash:   h_next,
  state_version:     prev+1,
  manifest_hash:     unchanged or new,
  content_hash:      unchanged or new,
  changed_fields:    ["manifest_hash"|"content_hash", ...],
  reason:            args.reason
}
\end{lstlisting}

\subsection{pages.archive@v1, pages.unpublish@v1, pages.delete@v1}

These three follow the same template: small \texttt{args}
($\texttt{slug}, \texttt{reason}$), a CanonicalReceipt with
correctly-set \texttt{pre\_state\_hash} / \texttt{post\_state\_hash}
/ \texttt{state\_version}. \texttt{pages.delete@v1} additionally
clears \texttt{content\_hash} and \emph{decrements} the blob's
ref-count (Phase 1 ref-count is a stub; see RFC-007).

\subsection{pages.get@v1 (query)}

\begin{lstlisting}
args:
{
  slug:    string,
  version: uint64|null      // null = current
}
returns:
{
  state_key:    ...,
  state_value:  string,
  canonical:    canonical_fields,
  body:         bytes|null,           // null if Deleted/Archived
                                      //   when caller is unauthorised;
                                      // omitted entirely if visibility
                                      //   gates the caller out
  version:      uint64,
  canonical_state_hash: bytes(32),
  etag:         bytes(16)             // §5.4
}
\end{lstlisting}

\texttt{pages.get@v1} produces an \texttt{OperationalReceipt}
recording the read, the caller URI, and (when from a Hub) the
\texttt{caller\_context} (main \S3.6).

\subsection{pages.list@v1 (query)}

\begin{lstlisting}
args:
{
  owner:        string,
  state_filter: ["Draft"|"Published"|...] | null,
  limit:        uint <= 1000,
  cursor:       string|null
}
returns:
{
  pages:        [{slug, state_value, version, etag, last_modified}],
  next_cursor:  string|null
}
\end{lstlisting}

% ──────────────────────────────────────────────────────────────────────────
\section{Algorithms}
\label{sec:algorithms}

\subsection{Canonical projection $\pi$}

\normsv{} Per main \S2.1, $\pi$ is schema-driven, not hand-rolled.
The pages schema lists exactly five fields with
\texttt{canonical:true}: \texttt{manifest\_hash},
\texttt{content\_hash}, \texttt{state\_value}, \texttt{visibility},
\texttt{version}. $\pi(v)$ projects $v$ onto exactly those keys, in
schema-declared order.

\begin{lstlisting}
fn project_pages(state: &PageState) -> Map<String, Value> {
    let mut out = BTreeMap::new();
    out.insert("manifest_hash".into(), bytes_to_hex(&state.manifest_hash));
    out.insert("content_hash".into(),  bytes_to_hex(&state.content_hash));
    out.insert("state_value".into(),   state.state_value.clone().into());
    out.insert("visibility".into(),    state.visibility.clone().into());
    out.insert("version".into(),       state.version.into());
    out  // BTreeMap iterates in lexicographic order — deterministic
}
\end{lstlisting}

\textbf{Cost.} $O(1)$ time and space (five fields, fixed sizes).

\subsection{Canonical serialisation $\sigma$ (JCS)}

\normsv{} $\sigma$ is RFC 8785 (JSON Canonicalization Scheme).
For pages this is a straight subset: object keys lexicographically
sorted (already guaranteed by \texttt{BTreeMap}), no extraneous
whitespace, hex strings for byte fields, integers without leading
zeros, no floats. JCS for the pages projection is therefore a
deterministic compaction of the projected JSON.

\subsection{canonical\_state\_hash $h$}

\normsv{} $h(b) = \mathrm{SHA256}(b)$, fixed in Phase 1.

\subsection{Worked byte example}

A page in Published state with the following canonical projection:

\begin{lstlisting}
{
  "content_hash":  "abf0c1...c8f3",   // 32 bytes hex
  "manifest_hash": "1d4a7e...09e2",
  "state_value":   "Published",
  "version":       3,
  "visibility":    "public"
}
\end{lstlisting}

After JCS:

\begin{lstlisting}
{"content_hash":"abf0c1...c8f3","manifest_hash":"1d4a7e...09e2","state_value":"Published","version":3,"visibility":"public"}
\end{lstlisting}

\noindent
$\texttt{canonical\_state\_hash} =
\mathrm{SHA256}(\text{above bytes})$
$= \texttt{0xa3b91f7e...8d04}$
(32 bytes). The first 16 bytes serve as the ETag
(\S\ref{sec:algo-concurrency}).

\subsection{Optimistic concurrency: ETag}
\label{sec:algo-concurrency}

\normsv{} Per RC-INV-2 (main \S3.7) every canonical transition
attempt declares its expected \texttt{pre\_state\_hash}. The Hub
and HTTP clients work better with a 16-byte ETag than a 32-byte
hex string; we therefore define:

\[
    \texttt{ETag}(p) = \mathrm{Base64Url}\bigl(\texttt{canonical\_state\_hash}(p)[0..16]\bigr)
\]

The 128-bit prefix is collision-safe at the population sizes a
single-realm Hub serves
(birthday bound $\approx 2^{64}$ pages before any collision).
On every \texttt{pages.update}, the caller supplies an
\texttt{If-Match} header (or, in the EasyNet ability path, an
\texttt{expected\_etag} field). The receipt-validation layer:

\begin{enumerate}[leftmargin=2.2em,itemsep=0.2em]
    \item Computes the page's current \texttt{ETag}.
    \item Compares against \texttt{expected\_etag}.
    \item On mismatch: rejects the transition with
    \texttt{ConcurrencyConflict}, recording the conflict in an
    \texttt{OperationalReceipt} (so the audit log records what
    was attempted but not committed).
\end{enumerate}

\begin{notebox}
\texttt{ETag} is a presentation of \texttt{canonical\_state\_hash}, not a
substitute. Internal protocol logic uses the full 32-byte hash
(RC-INV-2 cites \texttt{pre\_state\_hash}); the 16-byte ETag exists
only at the external HTTP boundary.
\end{notebox}

\subsection{View projection algorithms}
\label{sec:algo-views}

\subsubsection{human\_view (HTML)}

\normsv{}

\begin{lstlisting}
fn human_view(p: &PageState, blob: &BlobStore) -> HttpResponse {
    if p.state_value != "Published" {
        return http_404();
    }
    if !visibility_allows(p.visibility, request.caller) {
        return http_404();   // never 403 — leaks existence
    }
    let body = blob.get(&p.content_hash)?;
    let csp = "default-src 'self'; script-src 'none'; \
               style-src 'unsafe-inline' 'self'; img-src 'self' data:; \
               object-src 'none'; base-uri 'none'; frame-ancestors 'none'";
    HttpResponse::ok()
        .content_type(p.manifest.content_type)
        .header("ETag", etag(p))
        .header("Cache-Control", "public, max-age=60, must-revalidate")
        .header("Content-Security-Policy", csp)
        .header("Referrer-Policy", "no-referrer")
        .header("X-Content-Type-Options", "nosniff")
        .body(body)
}
\end{lstlisting}

\textbf{Determinism.} Per main TR-INV-6, \texttt{human\_view} is a
deterministic function of $(p, \text{blob})$. The CSP, the cache
header, and the response code are pure functions of \texttt{p}
and the request's principal; nothing here depends on wall-clock
time or load.

\textbf{Why 404 (not 403) on visibility miss.} A 403 confirms the
page exists. For \texttt{visibility=private} this is a leak. The
trade-off (private page indistinguishable from non-existent) is
intentional.

\subsubsection{agent\_view (structured)}

\normsv{}

\begin{lstlisting}
fn agent_view(p: &PageState) -> Value {
    let surface = derive_ability_surface(p);  // VP-INV-1
    let history = receipt_history_pointer(p); // §6.3
    json!({
        "state_type":  "easynet.page",
        "state_key":   { "realm": p.realm,
                         "owner": p.owner,
                         "slug":  p.slug },
        "ura":          p.ura(),
        "canonical": {
            "state":        p.state_value,
            "manifest":     p.manifest,         // small, in-line OK
            "content_hash": hex(p.content_hash),
            "visibility":   p.visibility,
            "version":      p.version,
            "etag":         etag(p)
        },
        "ability_surface":  surface,
        "history":  {
            "versions":                p.version,
            "latest_canonical_receipt_id": p.last_receipt_id,
            "receipt_log_ura":         history
        },
        "eal_hint": eal_hint_for(p),  // §5.6.3
        "view_meta": {
            "generated_from": "(canonical, runtime=∅, caller_context)"
        }
    })
}

fn derive_ability_surface(p: &PageState) -> Vec<&'static str> {
    match p.state_value.as_str() {
        "Draft"      => vec!["pages.publish@v1", "pages.update@v1",
                             "pages.delete@v1"],
        "Published"  => vec!["pages.update@v1", "pages.archive@v1",
                             "pages.delete@v1"],
        "Archived"   => vec!["pages.unpublish@v1", "pages.delete@v1"],
        "Deleted"    => vec![],
        _            => vec![],
    }
}
\end{lstlisting}

\textbf{Compliance with VP-INV-1.} \texttt{ability\_surface} is
derived strictly from \texttt{p.state\_value} and the static
transition schema (the match block). It is not a snapshot of the
ability registry; an ability that is registered but whose
\texttt{pre\_state} does not include \texttt{p.state\_value} is
invisible here.

\subsubsection{eal\_hint}

\normsv{} The \texttt{eal\_hint} is a single-line EAL fragment that
a peer agent can copy-paste to advance the page from its current
state. Examples:

\begin{lstlisting}
Draft     -> "call pages.publish on <owner> with { slug: \"<s>\" } timeout 30"
Published -> "call pages.update on <owner> with { slug: \"<s>\", body: <new> } timeout 60"
Archived  -> "call pages.unpublish on <owner> with { slug: \"<s>\" } timeout 30"
Deleted   -> ""    // no EAL hint; the surface is empty
\end{lstlisting}

The hint generator is deterministic; an agent reading the URA
twice without intervening transitions sees byte-identical hints.

% ──────────────────────────────────────────────────────────────────────────
\section{Architecture}
\label{sec:architecture}

\subsection{Module map (EasyNet-Cli)}
\label{sec:module-map}

\begin{lstlisting}
src/runtime/state/                   ← NEW (P-A1, P-A2)
  mod.rs                             pub use of types below
  store.rs                           StateStore trait + sled-backed impl
  hash.rs                            project / canonicalize_jcs / sha256
  receipt.rs                         CanonicalReceipt, OperationalReceipt

src/runtime/blob_store/              ← NEW (P-A1)
  mod.rs                             BlobStore trait
  fs.rs                              FsBlobStore (content-addressed
                                     directory tree, sha256 keys)

src/runtime/agents/pages_ability.rs  ← NEW (P-A2)
                                     register: create/publish/update/
                                     archive/unpublish/delete/get/list

src/runtime/views/                   ← NEW (P-A3)
  mod.rs                             View trait
  human.rs                           human_view (CSP, blob fetch)
  agent.rs                           agent_view (JSON projection)

src/bin/easynet-hub.rs               ← NEW (P-A4)
                                     axum binary, path → URA → view
src/runtime/hub/                     ← NEW (P-A4)
  mod.rs
  router.rs                          path parsing, view selection
  caller.rs                          principal/human-anon URI
                                     minting + caller_context
                                     extraction

src/runtime/eal_hint.rs              ← NEW (P-A5)
                                     deterministic EAL templates
\end{lstlisting}

\textbf{Design choice: \texttt{sled} for state and receipt index.}
sled is already a dependency in adjacent code; it gives us
range queries (needed for \texttt{state\_key}-prefix scans
mandated by main \S7.1) and atomic batches (needed to write a
canonical receipt + state mutation in one go). RFC-006 main does
not pin a backend; if sled is later replaced (e.g. by SQLite),
the \texttt{StateStore} trait stays.

\textbf{Design choice: filesystem blob store.} Phase 1 stores
blobs at \texttt{<state-root>/blobs/<aa>/<bb>/<full-hash>}, where
\texttt{aa} and \texttt{bb} are the first two hex bytes of the
SHA-256. The two-level fan-out keeps any single directory
under a few thousand entries even at large scale. RFC-007 will
refine this with ref-counts and a GC sweep; Phase 1 stubs both
(a blob is never deleted in Phase 1).

\subsection{Daemon wiring}

The daemon at boot:

\begin{enumerate}[leftmargin=2.2em,itemsep=0.2em]
    \item Opens the state store at \texttt{\$EASYNET\_HOME/state.sled}.
    \item Opens the blob store at \texttt{\$EASYNET\_HOME/blobs/}.
    \item Calls \texttt{pages\_ability::register} for every agent
    in the registry; the registration captures
    \texttt{(state\_store, blob\_store, agent\_name)} into the
    handler closure.
    \item Boots the receipt index (state\_key prefix index built
    on top of the same sled tree).
    \item Resumes any in-flight transitions
    (none in pages: every transition is single-shot, no
    \texttt{Building}-style intermediate state).
\end{enumerate}

\subsection{Hub wiring}

A separate binary, \texttt{easynet-hub}, embeds an axum router:

\begin{lstlisting}
GET  /<realm>/<owner>/<slug>           -> human_view
GET  /<realm>/<owner>/<slug>/.agent    -> agent_view (JSON)
GET  /<realm>/<owner>/<slug>/.receipts -> receipt history (JSON)
GET  /healthz                          -> 200 "ok"
\end{lstlisting}

The Hub is a daemon client: it has no direct access to the
state store. Every view fetch is a \texttt{pages.get@v1}
invocation against the daemon, with the envelope filled per
main \S5.4 (caller URI \texttt{principal/human-anon/...},
caller\_context filled per main \S3.6).

\textbf{Trust boundary.} The Hub is allowed to mint
\texttt{principal/human-anon} callers but not arbitrary
\texttt{agent/} callers. The daemon validates the caller URI on
admission and rejects \texttt{agent/} URIs that do not match
the Hub's authorisation token (an out-of-band shared secret in
Phase 1; RFC-008 will replace with proper signing).

\subsection{Storage layout}
\label{sec:storage}

\begin{lstlisting}
$EASYNET_HOME/
├── state.sled/                  sled DB
│   trees:
│     "page::canonical"          key: jcs((realm,owner,slug))
│                                value: bincode(PageState)
│     "page::receipts"           key: (jcs((realm,owner,slug)),
│                                      version_be_u64,
│                                      receipt_id)
│                                value: bincode(CanonicalReceipt)
│     "page::idx_owner"          key: (owner, slug)
│                                value: state_key (back-pointer)
│     "page::idx_receipt_id"     key: receipt_id
│                                value: state_key
└── blobs/
    ├── ab/
    │   ├── f0/
    │   │   └── abf0c1...c8f3    raw bytes
    └── ...
\end{lstlisting}

The two indexes (\texttt{idx\_owner}, \texttt{idx\_receipt\_id})
satisfy main \S7.1's dual-axis enumeration requirement: receipt
log enumerable both by \texttt{receipt\_id} (preserve original
EasyNet invocation traceability) and by \texttt{state\_key}
(new object-granularity traceability).

% ──────────────────────────────────────────────────────────────────────────
\section{Implementation plan}
\label{sec:impl-plan}

Five PR slices. Each slice is independently reviewable and
verifiable; merging the first does not block subsequent work in
parallel branches.

\subsection{P-A1: BlobStore + StateStore skeletons}

\textbf{Files added.} \texttt{src/runtime/blob\_store/} (\texttt{mod.rs}
+ \texttt{fs.rs}); \texttt{src/runtime/state/} (\texttt{mod.rs} +
\texttt{store.rs} + \texttt{hash.rs} + \texttt{receipt.rs}).

\textbf{Public API.}
\begin{lstlisting}
pub trait BlobStore: Send + Sync {
    fn put(&self, bytes: &[u8]) -> Result<[u8; 32]>;       // returns sha256
    fn get(&self, hash: &[u8; 32]) -> Result<Vec<u8>>;
    fn exists(&self, hash: &[u8; 32]) -> Result<bool>;
}

pub trait StateStore<S: StateObject>: Send + Sync {
    fn read(&self, key: &S::Key) -> Result<Option<S>>;
    fn apply(
        &self,
        receipt: &CanonicalReceipt,
        post_state: &S,
    ) -> Result<()>;     // atomic: state mutation + receipt append
    fn list_by_owner(...) -> Result<Vec<S::Key>>;
    fn enumerate_receipts(&self, key: &S::Key)
        -> Result<Vec<CanonicalReceipt>>;
}

pub fn project(state: &dyn StateObject) -> BTreeMap<String, Value>;
pub fn canonicalize_jcs(map: &BTreeMap<String, Value>) -> Vec<u8>;
pub fn canonical_state_hash(state: &dyn StateObject) -> [u8; 32];
\end{lstlisting}

\textbf{Tests.} JCS conformance (a curated set of RFC 8785 test
vectors); SHA-256 round-trip; \texttt{StateStore::apply}
atomicity under concurrent attempts (RC-INV-2 invariants pinned).

\textbf{Acceptance.} \texttt{cargo test -p easynet --lib state}
green; benchmarks $\geq$ 50k canonical hashes/sec on commodity
hardware (sanity bound, not a strict requirement).

\subsection{P-A2: pages.* ability registration + receipts}

\textbf{Files added.} \texttt{src/runtime/agents/pages\_ability.rs}.
Touches \texttt{src/runtime/agents/mod.rs} to register the
namespace.

\textbf{Behaviour.} Every transition listed in
\S\ref{sec:wire-schemas} is implemented:
\begin{enumerate}[leftmargin=2.2em,itemsep=0.15em]
    \item Validate args against the ability schema.
    \item Read current state via \texttt{StateStore}.
    \item Validate \texttt{pre\_state\_hash} (RC-INV-2).
    \item Compute \texttt{post\_state} and \texttt{post\_state\_hash}.
    \item Build \texttt{CanonicalReceipt}; sign with owner key
    (Phase 1: stub Ed25519 signature with a hard-coded keypair
    per agent; key management deferred to RFC-008).
    \item \texttt{StateStore::apply} (atomic).
    \item Return result.
\end{enumerate}

\textbf{Tests.}
\begin{itemize}[leftmargin=2.2em,itemsep=0.1em]
    \item Round-trip every transition.
    \item Reject \texttt{publish} from \texttt{Archived} (must
    \texttt{unpublish} first).
    \item Reject \texttt{create} with malformed slug.
    \item Reject any transition on \texttt{Deleted} (terminal).
    \item Optimistic concurrency: two \texttt{update} attempts
    racing on the same page \texttt{slug}; one wins, one rejected
    with \texttt{ConcurrencyConflict}.
\end{itemize}

\textbf{Acceptance.} The ability is callable via
\texttt{easynet ability invoke}; receipts visible via the
existing receipt log surface.

\subsection{P-A3: view projection + EAL hints}

\textbf{Files added.} \texttt{src/runtime/views/} +
\texttt{src/runtime/eal\_hint.rs}.

\textbf{Behaviour.} Pure functions: a \texttt{PageState} maps to
a \texttt{(human\_view\_response, agent\_view\_json)} pair plus
an EAL hint string. No network, no state mutation.

\textbf{Tests.} Determinism (1000 iterations of identical input
yield identical output bytes); \texttt{ability\_surface} matches
\S\ref{sec:state-space} for every state; CSP header constant.

\subsection{P-A4: Hub binary}

\textbf{Files added.} \texttt{src/bin/easynet-hub.rs} +
\texttt{src/runtime/hub/}.

\textbf{Behaviour.} axum router, three GET routes, each:
\begin{enumerate}[leftmargin=2.2em,itemsep=0.1em]
    \item Parse path $\to$ \texttt{(realm, owner, slug)} $\to$
    \texttt{state\_key}.
    \item Build envelope: caller URI is
    \texttt{principal/human-anon/<ip-hash>/<session-id>}
    (\S\ref{sec:hub}); \texttt{caller\_context} populated from
    HTTP headers.
    \item Invoke \texttt{pages.get@v1} on the daemon.
    \item Apply view projection (P-A3).
    \item Return HTTP response.
\end{enumerate}

\textbf{Tests.} Integration: launch a daemon + Hub in a temp dir,
\texttt{create+publish} a page, \texttt{curl} through the Hub,
assert HTML response + correct ETag + correct CSP. \texttt{curl}
the \texttt{/.agent} suffix, assert JSON shape matches
\S\ref{sec:algo-views}.

\subsection{P-A5: agent\_view enrichment + receipt history endpoint}

\textbf{Behaviour.} Implements the \texttt{/.receipts} route:
returns the receipt history for a page (canonical receipts only,
in version order) as a JSON array. Adds an
\texttt{eal\_hint} field to \texttt{agent\_view} (already
specified in P-A3 but only stubbed there).

\textbf{Tests.} Replay a receipt history client-side using
\texttt{canonical\_state\_hash} reconstruction; assert the final
hash equals the page's current hash.

\subsection{Hub envelope minting}
\label{sec:hub}

\begin{lstlisting}
fn build_envelope(req: &HttpRequest) -> Envelope {
    let ip   = client_ip(req);
    let salt = HUB_LOCAL_SALT.read();
    let ip_hash = hmac_sha256(&salt, ip.to_string().as_bytes())
        .iter().take(8)
        .flat_map(|b| format!("{:02x}", b).chars().collect::<Vec<_>>())
        .collect::<String>();
    let sid = req.cookie("easynet_session")
        .or_else(|| req.header("X-EasyNet-Session"))
        .or_else(|| req.query("session"))
        .unwrap_or_else(|| "-".to_string());

    let caller = format!(
        "easynet:///principal/human-anon/{ip_hash}/{sid}"
    );
    let caller_context = json!({
        "origin": {
            "referer_url":   req.header("Referer"),
            "referer_realm": referer_origin(req.header("Referer")),
            "direct":        is_direct_navigation(req)
        },
        "agent_class": classify_user_agent(req.header("User-Agent")),
        "headers": {
            "accept_language":      req.header("Accept-Language"),
            "sec_fetch_site":       req.header("Sec-Fetch-Site"),
            "sec_fetch_mode":       req.header("Sec-Fetch-Mode"),
            "sec_fetch_dest":       req.header("Sec-Fetch-Dest"),
            "x_forwarded_for_hash": xff_hash(req)
        }
    });
    Envelope { caller, caller_context, ... }
}
\end{lstlisting}

% ──────────────────────────────────────────────────────────────────────────
\section{End-to-end worked example}
\label{sec:e2e}

A complete trace from \texttt{pages.create} through Hub-rendered
HTML access and agent URA fetch.

\subsection{Setup}

\begin{lstlisting}
$ easynet agent add alice --type claude-code
$ easynet-daemon &       # listens on $EASYNET_HOME ipc
$ easynet-hub --bind 127.0.0.1:8787 &
\end{lstlisting}

\subsection{Author and publish}

\begin{lstlisting}
$ easynet ability invoke alice.pages.create --args '{
    "slug": "shop-001",
    "manifest": {
        "title": "Mug, ceramic, 12oz",
        "content_type": "text/html",
        "visibility": "public"
    },
    "initial_body": "<h1>Mug</h1><p>$19.99 — ships Tuesday.</p>"
}'

response:
  receipt_id:           "r_a1b2..."
  state_value:          "Draft"
  canonical_state_hash: "0xa3b91f7e..."
  version:              1

$ easynet ability invoke alice.pages.publish --args '{ "slug": "shop-001" }'

response:
  receipt_id:           "r_c3d4..."
  state_value:          "Published"
  canonical_state_hash: "0x4d28ee17..."
  version:              2
  etag:                 "TSjuFw0pAvjpAQ=="
\end{lstlisting}

\subsection{Browser fetch via Hub}

\begin{lstlisting}
$ curl -i http://127.0.0.1:8787/default/alice/shop-001

HTTP/1.1 200 OK
ETag: "TSjuFw0pAvjpAQ=="
Cache-Control: public, max-age=60, must-revalidate
Content-Security-Policy: default-src 'self'; script-src 'none'; ...
Content-Type: text/html
Content-Length: 49

<h1>Mug</h1><p>$19.99 — ships Tuesday.</p>
\end{lstlisting}

\noindent The Hub emits an \texttt{OperationalReceipt} for the
read:

\begin{lstlisting}
{
  "receipt_class":  "operational",
  "state_key":      ["default","alice","shop-001"],
  "transition_id":  "pages.get@v1",
  "attempt_id":     "att_e5f6...",
  "envelope": {
    "caller": "easynet:///principal/human-anon/9c4e.../sess-abc",
    "caller_context": {
      "origin":      { "referer_url":   null,
                       "referer_realm": null,
                       "direct":        true },
      "agent_class": { "kind":          "browser",
                       "user_agent_hash": "a1b2..." },
      "headers":     { "accept_language": "en-US",
                       "sec_fetch_site":  "none",
                       "sec_fetch_mode":  "navigate",
                       "sec_fetch_dest":  "document" }
    }
  },
  "observed_at": "2026-04-28T10:14:32Z"
}
\end{lstlisting}

\subsection{Peer agent fetches via URA}

\begin{lstlisting}
$ curl http://127.0.0.1:8787/default/alice/shop-001/.agent | jq

{
  "state_type":  "easynet.page",
  "state_key":   { "realm": "default",
                   "owner": "alice",
                   "slug":  "shop-001" },
  "ura":         "easynet:///r/default/agent/alice/page/shop-001",
  "canonical": {
    "state":        "Published",
    "manifest":     { "title":         "Mug, ceramic, 12oz",
                      "content_type":  "text/html",
                      "visibility":    "public" },
    "content_hash": "0x4d28ee17...",
    "visibility":   "public",
    "version":      2,
    "etag":         "TSjuFw0pAvjpAQ=="
  },
  "ability_surface": ["pages.update@v1", "pages.archive@v1",
                      "pages.delete@v1"],
  "history": {
    "versions": 2,
    "latest_canonical_receipt_id": "r_c3d4...",
    "receipt_log_ura":
      "easynet:///r/default/agent/alice/page/shop-001/receipts"
  },
  "eal_hint":
    "call pages.update on alice with { slug: \"shop-001\", body: <new> } timeout 60"
}
\end{lstlisting}

\subsection{Audit replay}

A peer fetches the receipt history and replays it locally:

\begin{lstlisting}
$ curl http://127.0.0.1:8787/default/alice/shop-001/.receipts | jq -r '.[].post_state_hash'
0xa3b91f7e...   # version 1, after create
0x4d28ee17...   # version 2, after publish
\end{lstlisting}

\noindent The peer:
\begin{enumerate}[leftmargin=2em,itemsep=0.15em]
    \item Reconstructs canonical state from each receipt
    in order.
    \item Computes \texttt{canonical\_state\_hash} for each
    reconstruction.
    \item Asserts each computed hash equals the receipt's
    \texttt{post\_state\_hash}.
    \item Asserts the final hash equals the agent\_view's
    \texttt{canonical.content\_hash} (well, not literally ---
    the comparison is on \texttt{canonical\_state\_hash}, but
    \texttt{content\_hash} is one of its inputs).
    \item Verifies each receipt's \texttt{owner\_signature}
    against \texttt{alice}'s published key.
\end{enumerate}

\noindent Successful audit means: the peer has cryptographic
proof that this page, at version 2, presented exactly this
content, signed by Alice, on or before the timestamp in the
final receipt. This is the audit story the traditional web
cannot tell.

% ──────────────────────────────────────────────────────────────────────────
\section{Open questions deferred to RFC-006-A v2}
\label{sec:open}

\begin{enumerate}[leftmargin=2.2em,itemsep=0.2em]
    \item \textbf{HTML sanitisation.} Phase 1 trusts owner-authored
    HTML entirely (CSP \texttt{script-src 'none'} mitigates the
    worst case but still lets \texttt{onerror=}-style attribute
    injection through). v2 should specify a safe-subset
    sanitiser run at \texttt{publish} time and pin its output
    hash.
    \item \textbf{Multi-realm Hubs.} A Hub serving more than one
    realm needs a routing table; the path scheme
    \texttt{/<realm>/...} already accommodates it but Phase 1
    binds a single realm at startup.
    \item \textbf{Page templates / slot-filling.} Common pages
    (product catalogue, blog post) repeat structure. A template
    layer could let owners author manifests + slot data and
    have the canonical \texttt{content\_hash} cover the
    rendered output.
    \item \textbf{Authenticated humans
    (\texttt{principal/human-auth}).} For private/realm pages
    that need real user identity, the
    \texttt{principal/human-auth/<provider>/<user-id>} sub-namespace
    must be specified along with its signature scheme.
    \item \textbf{Federated reads.} A page owned by node A read
    through a Hub on node B. The Hub becomes a federation
    client, not just a daemon client. RFC-006-A v2.
\end{enumerate}

\subsection{Coverage check}

\begin{center}
\renewcommand{\arraystretch}{1.2}
\begin{tabularx}{\linewidth}{l X l}
\toprule
\textbf{Main rule} & \textbf{Discharged by} & \textbf{Where} \\
\midrule
SO-INV-1 (key immutability)              & state\_key parser refuses key updates                 & \S\ref{sec:state-machine} \\
CSH-INV-1 (hash over projection)         & \texttt{project\_pages} + JCS only                     & \S\ref{sec:algorithms} \\
CSH-INV-2 (content as blob)              & FsBlobStore; canonical state holds only the hash       & \S\ref{sec:storage} \\
RC-INV-1 (receipt binding)               & every receipt carries (state\_key, transition\_id, attempt\_id) & \S\ref{sec:wire-schemas} \\
RC-INV-2 (optimistic concurrency)        & ETag check at receipt-validation time                  & \S\ref{sec:algo-concurrency} \\
TA-INV-1 (transition\_id required)       & every \texttt{pages.<verb>@v1} string is validated     & \S\ref{sec:wire-schemas} \\
VP-INV-1 (ability\_surface from state)   & \texttt{derive\_ability\_surface} match block          & \S\ref{sec:algo-views} \\
TR-INV-7 (multi-view first-class)        & \texttt{human\_view} + \texttt{agent\_view} both implemented & \S\ref{sec:algo-views} \\
TR-INV-11 (URA-receipt traceability)     & \texttt{/.receipts} endpoint + idx\_owner index        & \S\ref{sec:storage}, \S\ref{sec:e2e} \\
TR-INV-12 (principal/human-anon)         & \texttt{build\_envelope} mints the URI                 & \S\ref{sec:hub} \\
TR-INV-13 (caller\_context)              & populated by \texttt{build\_envelope}; entered into operational receipts & \S\ref{sec:hub}, \S\ref{sec:e2e} \\
\bottomrule
\end{tabularx}
\end{center}

\paragraph{Closing.} Pages is the simplest reference machine that
exercises the full normative core of stateful EasyNet.
Implementing the five PRs above closes Phase 1 of stateful
EasyNet and unlocks both the pressure-test machine
(\texttt{webapp}, Appendix B) and the agent-native web layer
that follows.

\vfill
\hrule
\vspace{0.4em}
\noindent\small
\textit{This document is a Phase-1 implementation specification.
Discrepancies with AXON-RFC-006 main are bugs in this document;
discrepancies with Appendix~B's pressure tests are bugs in either
this document or the main body, to be reconciled in alignment.}

\end{document}
